\section{结束语}

本文围绕"生态格局—分析骨架—路线比较—关键挑战—收敛方向"这一主线，对国产算力平台上的大模型推理引擎作了结构化梳理。全文的核心判断可归纳为：国产推理引擎的竞争点不在局部算子是否足够快，而在于系统能否在同一架构中同时承受后端异构、复杂任务负载和上游快速演进三类压力，并在国产芯片与本土部署环境中形成可持续的工程闭环。

具体而言，这一判断体现在四个方面。第一，平台异构的真正终点不是兼容，而是硬件差异能否被收敛为稳定、可验证、可回退的能力契约——若能力边界长期以隐式特判存在，新后端的接入就会反复冲击共享主链。第二，长上下文、多模态、结构化输出与 Agent 负载正在把推理引擎从 token 生成器推向状态中心化运行时，状态治理能力将比局部算子优化更能决定系统上限。第三，量化与 KV 压缩只有进入统一的观测、调度和成本闭环，才可能从离线实验结果转化为可持续交付的单位 token 成本优势。第四，路线选择本质上是复杂度分配与责任边界选择——复杂度由 backend、runtime、state machine 还是控制面承担，将直接决定系统的演进速度与长期维护成本。

面向未来，国产推理引擎的下一阶段竞争不在单点峰值还能提高多少，而在于谁能更快、更稳、更低成本地吸收模型变化并完成生产交付。对产业侧而言，建设重点应从围绕单模型做专项优化转向围绕持续演进负载建设系统能力；对研究侧而言，则需要把平台迁移成本、长周期维护成本和近真实服务指标放到与性能同等重要的位置。围绕这些问题持续推进，国产推理引擎的演进才更可能从阶段性适配走向可持续积累。
